% !TeX root = ../../Book1.tex
\section{定向、体积形式与参数选择}
\label{b1:sec:B105}

面积公式中的 Jacobi 因子总是非负数，而通常行列式可以为负。负号并不是需要修正的面积误差；它记录了映射是否反转方向。对标量函数作曲面积分时方向可以忽略，对交错型积分或讨论曲面的边界时却不能。本节先在一个向量空间内固定符号规则，再将它用于连续变化的切空间。

\subsection{顶次外空间中的两个方向}
\label{b1:subsec:B10501}

\begin{definition}\label{b1:def:app-b-orientation}
设实向量空间 \(V\) 的维数为 \(n\geq1\)。对两组有序基 \(e=(e_1,\ldots,e_n)\) 与 \(f=(f_1,\ldots,f_n)\)，若由
\[
 f_j=\sum_{i=1}^n a_{ij}e_i
\]
定义的过渡矩阵满足 \(\det A>0\)，则称这两组基同向。选择全部有序基在此关系下的一个等价类，称为给 \(V\) 指定一个\term[dingxiang]{定向}[orientation]；被选类中的基称为正向基，另一类中的基称为负向基。
\end{definition}

行列式的乘法公式说明同向关系满足自反性、对称性与传递性。固定一组基后，其他基的过渡矩阵行列式只能为正或为负；交换固定基的前两个向量，或把第一个向量取负，可以互换这两种符号。因此 \(V\) 恰有两个定向。这个选择不依赖长度，也不需要内积。

\ProseParagraphBreak

外代数把这一定义压缩到一个一维空间中。由于
\[
 f_1\wedge\cdots\wedge f_n
  =(\det A)e_1\wedge\cdots\wedge e_n,
\]
同向的基恰好给出 \(\bigwedge^nV\setminus\{0\}\) 中同一条开射线。指定定向等于从这条实直线的两个方向中选一个方向，而不是选定某个特定长度的向量。

\ProseParagraphBreak

对任意线性自映射 \(A:V\rightarrow V\)，顶次外映射作用在一维空间上，必是乘以某个标量。基下的子式公式给出
\[
 \bigwedge\nolimits^n A=(\det A)I_{\bigwedge^nV}.
\]
因此行列式可以不借助具体矩阵而定义；换基后标量不变，乘法公式来自外映射与复合相容。若 \(A\) 可逆，正行列式对应保持定向，负行列式对应反转定向。奇异映射把顶次外空间压为零，此时不能谈它把一个定向送到了像的哪个 \(n\) 维定向。

\ProseParagraphBreak

给 \(V\) 增加欧氏内积后，每个定向都有唯一的正向单位顶次外向量
\[
 \tau_V=e_1\wedge\cdots\wedge e_n,
\]
其中 \(e_i\) 为任意正向标准正交基。另一组同向标准正交基的过渡矩阵行列式为 \(1\)，故 \(\tau_V\) 与基无关。反转定向只把 \(\tau_V\) 换成 \(-\tau_V\)。

\begin{proposition}\label{b1:prop:app-b-volume-form}
在已定向的 \(n\) 维欧氏空间 \(V\) 上，存在唯一的顶次交错型 \(\mathrm{vol}_V\)，使
\[
 \langle\mathrm{vol}_V,\tau_V\rangle=1.
\]
它称为与内积及定向相应的\term[tiji-xingshi]{体积形式}[volume form]。若 \(v_j=\sum_i a_{ij}e_i\)，其中 \(e_i\) 是正向标准正交基，则
\[
 \mathrm{vol}_V(v_1,\ldots,v_n)=\det A,\quad
 \bigl|\mathrm{vol}_V(v_1,\ldots,v_n)\bigr|
       =|v_1\wedge\cdots\wedge v_n|.
\]
若 \(V,W\) 都是已定向的 \(n\) 维欧氏空间，\(L:V\rightarrow W\) 线性，则
\[
 L^*\mathrm{vol}_W=(\det_{\mathrm{or}}L)\mathrm{vol}_V,
 \quad
 |\det_{\mathrm{or}}L|=J_n(L),
\]
其中 \(\det_{\mathrm{or}}L\) 是 \(L\) 在两端正向标准正交基下的行列式。
\end{proposition}
\begin{proof}
顶次外空间是一维的，所以在非零向量 \(\tau_V\) 上指定值便唯一决定一个线性泛函；再由外空间与交错型的配对识别得到 \(\mathrm{vol}_V\)。将
\(v_1\wedge\cdots\wedge v_n=(\det A)\tau_V\) 代入，就得到第一个公式。由于 \(|\tau_V|=1\)，取绝对值得到第二个公式。

最后，\(L^*\mathrm{vol}_W\) 仍是 \(V\) 上的顶次型，故为 \(\mathrm{vol}_V\) 的标量倍数。在正向标准正交基上求值，倍数恰为所写行列式。其绝对值由 Gram 公式等于面积因子。两端正向标准正交基的改变各自只有行列式 \(1\)，所以这个带符号的数不随所用基而变。
\end{proof}

欧氏体积测度与体积形式要区分。前者不给非负函数的积分增加符号；后者是作用于有序向量组的交错型，改变定向便取负。若把内积换掉而保持定向，正向基的集合不变，单位顶次外向量和体积形式的具体大小则会改变。

\ProseParagraphBreak

以下不是由前述有序基的等价类定义推出，而是本书为零维定向边界单独采用的符号约定：令其顶次外空间为 \(\bigwedge^0\{0\}\coloneqq\mathbb R\)，并用符号 \(+1\) 或 \(-1\) 指定定向。唯一的空基不能区分这两个选择。因而一个定向点带有正号或负号；有限多个定向点的积分就是带这些符号的求和。

\subsection{切空间的连续定向}
\label{b1:subsec:B10502}

设 \(M\subseteq\mathbb R^N\) 是 \(k\) 维 \(C^1\) 子流形。这里的局部描述是：每点附近都可由一个 \(C^1\) 参数化
\[
 f:U\subseteq\mathbb R^k\longrightarrow M
\]
给出，\(f\) 在该邻域是到其像的同胚，且 \(Df\) 处处具有秩 \(k\)。切空间为 \(T_{f(s)}M=Df(s)(\mathbb R^k)\)。本节只使用这类局部参数图及其 \(C^1\) 过渡映射；不在这里展开一般抽象流形理论。

\ProseParagraphBreak

参数域的标准定向给出局部连续的单位可分解外向量场
\[
 \tau_f\bigl(f(s)\bigr)
 =\frac{\partial_1f(s)\wedge\cdots\wedge\partial_kf(s)}
        {J_kf(s)}.
\]
分母处处为正，所以该式确实连续。若同一部分曲面另用 \(g=f\circ\phi\) 参数化，其中 \(\phi\) 为两个参数开集之间的 \(C^1\) 微分同胚，则
\begin{equation}\label{b1:eq:app-b-param-orientation}
 \tau_g\bigl(g(t)\bigr)
   =\operatorname{sgn}\bigl(\det D\phi(t)\bigr)
      \tau_f\bigl(g(t)\bigr).
\end{equation}
这是上一节外映射的复合公式与 \(J_k(g)=J_k(f\circ\phi)\) 的直接商式。过渡行列式连续且从不为零，所以其符号在每个连通分支上固定。

\ProseParagraphBreak

若能在 \(M\) 上连续选择单位可分解外向量 \(\tau(x)\)，使其所确定的平面为 \(T_xM\)，则称 \(M\) 可定向；选定这样的 \(\tau\) 就给出了 \(M\) 的定向。局部参数图若满足 \(\tau_f=\tau\)，便是正向参数图。由公式\eqref{b1:eq:app-b-param-orientation}，正向参数图之间的过渡行列式为正；反过来，若选出的参数图覆盖 \(M\)，且所有过渡映射的行列式均为正，那么这些局部 \(\tau_f\) 在重叠处相等，可以拼成全域连续的 \(\tau\)。

\ProseParagraphBreak

每张参数图都可定向，并不保证这些选择能在整张曲面上拼接。问题出在沿着不同图走一圈后符号能否保持一致，不出在某一点的线性代数。若 \(M\) 连通且已有一个定向 \(\tau\)，任何其他连续定向只能写成 \(\sigma(x)\tau(x)\)，其中 \(\sigma(x)\in\{-1,1\}\) 连续。连通性使 \(\sigma\) 为常数，故这时恰有两个全域定向。若有多个连通分支，各分支可以独立选择符号。

\ProseParagraphBreak

在\chapref{b1:ch:19}的可整流情形，切空间一般只在几乎处处存在；\appref{b1:app:D}将用可测的单位可分解外向量记录方向，而不要求它连续。连续定向与可测定向具有不同的用途，不能把流形的全局可定向性当作所有可整流积分对象的先决条件。

\subsection{外法向优先的边界约定}
\label{b1:subsec:B10503}

设 \(V\) 是已定向的 \(n\) 维欧氏空间，\(\nu\) 是一个单位向量，令 \(H=\nu^\perp\)。给 \(H\) 选择定向时，规定把 \(\nu\) 放在其正向基的最前面，所得到的 \(V\) 的基应为正向。这称为外法向优先约定；此处“外”是否具有几何含义，取决于 \(\nu\) 是否来自某个区域的外法向。

\begin{proposition}\label{b1:prop:app-b-boundary-orientation}
当 \(n\geq2\) 时，存在唯一的单位可分解外向量
\(\tau_H\in\bigwedge^{n-1}H\)，使
\[
 \nu\wedge\tau_H=\tau_V.
\]
它对应于上述定向。其体积形式由
\[
 \mathrm{vol}_H(v_1,\ldots,v_{n-1})
   =\mathrm{vol}_V(\nu,v_1,\ldots,v_{n-1}),
 \quad v_i\in H
\]
给出。反转 \(\nu\) 而保持 \(V\) 的定向，或反转 \(V\) 的定向而保持 \(\nu\)，都会反转 \(H\) 的定向。
\end{proposition}
\begin{proof}
选 \(H\) 的一组标准正交基 \(q_1,\ldots,q_{n-1}\)，则
\(\nu,q_1,\ldots,q_{n-1}\) 是 \(V\) 的标准正交基，故其外积为 \(\pm\tau_V\)。如有必要把 \(q_1\) 取负，就可令符号为正。由此构造了所需的 \(\tau_H\)。

空间 \(\bigwedge^{n-1}H\) 为一维，且映射 \(\eta\mapsto\nu\wedge\eta\) 把其非零基向量送到非零向量，因此单射，证明唯一性。右侧定义的型在这组正向标准正交切向量上取值为 \(1\)，所以按体积形式的唯一性就是 \(\mathrm{vol}_H\)。最后两种符号改变分别代入 \(\nu\wedge\tau_H=\tau_V\)，便得到相应的 \(\tau_H\) 都取负。
\end{proof}

若 \(\Omega\subseteq\mathbb R^n\) 有 \(C^1\) 边界，采用环境标准定向，并把连续单位外法向 \(\nu_\Omega\) 代入该命题，就得到边界定向。例如局部写
\[
 \Omega=\{x:F(x)<0\},\quad \nabla F\ne0\quad\text{于边界},
\]
则 \(\nu_\Omega=\nabla F/|\nabla F|\)。使用另一个表示同一内部一侧的定义函数，不会改变外法向。该约定与\chapref{b1:ch:27}的外法向一致；有限周长集合上的相应定向须在约化边界上作几乎处处的解释。

\ProseParagraphBreak

在平面单位圆上，外法向为
\(\nu(\theta)=(\cos\theta,\sin\theta)\)，而
\[
 \tau(\theta)=(-\sin\theta,\cos\theta)
\]
满足 \(\det(\nu,\tau)=1\)，所以正向是逆时针。若考虑一个圆环，内边界的区域外法向指向圆心，故内圆的正向是顺时针。两个边界分支的方向相反，来自同一个“区域的外侧”条件，不是额外补出的符号规则。

\ProseParagraphBreak

对图面 \(f(s)=(s,g(s))\subset\mathbb R^{k+1}\)，向上的单位法向为
\[
 \nu=\frac{(-\nabla g,1)}{\sqrt{1+|\nabla g|^2}}.
\]
列矩阵 \((\partial_1f,\ldots,\partial_kf,\nu)\) 的行列式为
\(\sqrt{1+|\nabla g|^2}>0\)。将最后一列移到第一列需要交换 \(k\) 次，因此外法向优先约定给出的切向外向量是
\[
 \tau_H=(-1)^k
 \frac{\partial_1f\wedge\cdots\wedge\partial_kf}
      {\sqrt{1+|\nabla g|^2}}.
\]
二维曲面的熟悉规则与一维曲线的规则因此看起来不同：关键是法向放在首位还是末位，不能在维数改变时省去这一次符号核对。

\ProseParagraphBreak

在 \(n=1\) 时把上述方程解释为 \(\nu\,\tau_H=\tau_V\)，其中 \(\tau_H\in\{-1,1\}\)。标准定向区间 \([a,b]\) 的左端外法向为 \(-1\)，右端为 \(1\)，所以其定向边界为“右端点减左端点”。这正与微积分基本定理中的 \(f(b)-f(a)\) 对应。

\subsection{定向积分与参数变换}
\label{b1:subsec:B10504}

现在设 \(\omega(x)\) 是定义在曲面附近的连续 \(k\) 次交错型场，即各点都有一个交错型，且其在固定坐标基上的系数连续。给曲面 \(M\) 指定单位切向外向量场 \(\tau\)，在
\[
 \int_M\bigl|\langle\omega(x),\tau(x)\rangle\bigr|
       \,\mathrm d\mathcal H^k(x)<\infty
\]
时，定义这个型场在定向曲面上的积分为
\[
 \int_{(M,\tau)}\omega
   \coloneqq\int_M\langle\omega(x),\tau(x)\rangle
          \,\mathrm d\mathcal H^k(x).
\]
右侧是本书已经建立的标量积分，所以定义先有明确的可积性条件，再讨论参数表达。这里不要求提前掌握微分形式的外微分或 Stokes 定理。

\ProseParagraphBreak

若 \(f:U\rightarrow M\) 是覆盖所考察曲面片的正向单射 \(C^1\) 参数图，则 \(f\) 在开集 \(U\) 上局部 Lipschitz。先把\cref{b1:cor:local-area-formulas}的非负加权局部面积公式用于
\(\bigl|\langle\omega(f(s)),\tau(f(s))\rangle\bigr|\)，由上面的可积性假设得到相应参数域积分有限；再使用同一结果的实值加权版本，便有
\[
 \begin{aligned}
 \int_{(M,\tau)}\omega
 &=\int_U
   \langle\omega\bigl(f(s)\bigr),\tau\bigl(f(s)\bigr)\rangle
       J_kf(s)\,\mathrm ds\\
 &=\int_U
   \omega\bigl(f(s)\bigr)(\partial_1f(s),\ldots,\partial_kf(s))
       \,\mathrm ds.
 \end{aligned}
\]
分母中的面积因子被抵消，剩下的是直接测试有序偏导向量。若改用负向参数图，同一个等式的最后一行要取负。

\ProseParagraphBreak

具体地，设 \(g=f\circ\phi\)，且 \(\det D\phi\) 在所用参数域上恒为负。型的多线性给出
\[
 \omega\bigl(g(t)\bigr)(\partial_1g(t),\ldots,\partial_kg(t))
 =(\det D\phi(t))\,
   \omega\Bigl(f\bigl(\phi(t)\bigr)\Bigr)
       \Bigl(\partial_1f\bigl(\phi(t)\bigr),\ldots,
             \partial_kf\bigl(\phi(t)\bigr)\Bigr).
\]
普通变量代换使用 \(|\det D\phi|\)，因而直接把 \(g\) 当作正向参数图来积分，就会得到原积分的负值。若保持曲面的同一个定向，则必须同时记入 \(g\) 的负向符号，两者抵消后积分不变。若有多个参数分支，分别在符号固定的分支上讨论即可。

\ProseParagraphBreak

这说明无向面积公式中的绝对值和交错型拉回中的带符号行列式并不矛盾。前者改变积分所用的非负测度，后者改变被积对象对有序向量的取值。将定向整体反转，面积测度不动而 \(\tau\) 取负，所以型的积分整体取负。\appref{b1:app:D}把这种积分看作作用于测试形式的线性泛函，边界也将在这一对偶语言中统一处理。
